\section{Methodology} \label{sec:methodology}
\subsection{Existing approaches and related literature}

The measurement of the group based mean wage differences can be traced back to the seminal work of \textcite{Oaxaca1973} and \textcite{Blinder1973}, wherein the mean wage gap was decomposed into composition and structure effects. This methodological development has spurred a substantial body of research aimed at refining and extending the Oaxaca-Blinder decomposition beyond the single point estimate of the wage distribution over the past half-century \parencite{Fortin2011}. Methodological extensions towards whole of the wage distribution allows for the identification of gaps specific to particular wage groups, facilitating a deeper understanding of differences both between and within groups \parencite{Machado2005}. Majority of literature on wage gap beyond the mean focuses on Conditional Quantile Regression (CQR) models \parencite{Machado2005,Buchinsky1998}; however, the method is path dependent\footnote{The decomposition approach depends on the order in which the decomposition is performed \textcite{Fortin2011}}.  In contrast, we use Unconditional Quantile Regressions (UQR)\footnote{Using the CQRs,  we can only interpret the effect of a unit change in a covariate X on the t-th quantile of the conditional quantile distribution whereas, UQRs allows to estimate the marginal effects of explanatory variables on the unconditional distribution of y and therefore on distributional statistics like quantiles.} of RIFs based on  {Nopo2009} to obtain a Oaxaca-Blinder based detailed decomposition where linearity is assumed and is path independent. UQR model over CQR model allows for the unconditional mean interpretation and enables to solve the problem of RIF-OLS model misspecification due to nonlinearity with a semiparametric reweighting scheme based on \textcite{DiNardo1996}.  \par

A major hurdle in distributional decomposition is to construct a counterfactual distribution, which can not be directly observed. As a result, significant amount of effort in decomposition literature has been devoted to develop methods for constructing counterfactual. \textcite{DiNardo1996} use kernel density reweighting, whereas, \textcite{Firpo2009} utilize recentered influence function. \textcite{Machado2005} deploy quantile regression to estimate the inverse conditional distribution function. In contrast, \textcite{Chernozhukov2013} tackle the problem directly by estimating the conditional distributional regression model using quantile regression.\par

In addition to going beyond mean, addressing the selection concern has been an important agenda in studying gender wage differentials. Four major strategies in the literature have been developed, namely: (a) imputation, (b) identification at infinity, (c) parametric modeling of selection, and (d) the bounding approach \parencite{Machado2017}. The imputation method involves utilizing observed covariates and economic model-based restrictions to impute values for the missing part of the data, i.e., those who do not participate in the work. A recent application, \textcite{Blau2021}, searches backward and forward in the panel data and proxies missing wage by the observation in the nearest wave. In contrast, identification at infinity circumvents the selection by limiting itself only to a much-smaller segment of labor force where participation rates are very high and selection is considered negligible \parencite{Machado2017,Mulligan2008,Heckman1990}.\par

The parametric approach to selection correction is bolder; it aims to explicitly model the selection process, either at the mean \parencite{Newey2009, Heckman1979, Heckman1974} or at quantiles \parencite{Buchinsky1998}. In these models, the outcome and the latent selection equations exhibit linearity with respect to covariates and error terms are assumed to be independent of covariates conditional on the selection probability. In comparison, the bounding approach has a lesser ambition as it only seeks to tighten the worst-case scenario bounds on the gender wage gap vi\'{a} restrictions motivated by the economic theory \parencite{Blundell2007}. But, these restrictions -- availability of instrument to tighten the bound, pre-suppositions on the selection's sign, or both -- being weaker than parametric modeling impose wider bounds.\par

In the spirit of \textcite{Buchinsky1998}, we correct for selection vi\'{a} parametric approach in  the quantile framework. However, we use \textcite{Arellano2017}'s copula based technique to model the joint-distribution of error terms in outcome and selection models. This approach overcomes \textcite{Huber2015}'s  critique concerning the conditional independence assumption in sample selection models, particularly its implication of identical slopes across all quantile regressions. With additional restrictions compared to the bounding approach, our methodology provides more tighter bounds and greater flexibility in capturing the direction of sample selection from the observed data, rather than relying solely on theoretical priors.\par


\subsection{RIF-regression}
The influence function (IF) is defined as IF($y; v$), the influence function corresponding to an observed wage $y$ for the distributional statistic of interest, $v(F_Y)$. The recentered influence function (RIF) is defined as RIF($y; v$) = $v(F_Y) + \text{IF}(y; v)$, so that it aggregates back to the statistics of interest ($\mathbb{E}[\text{RIF}(y; v)] = v(F_Y)$). In its simplest form, the approach assumes that the conditional expectation of the RIF($y; v$) can be modeled as a linear function of the explanatory variables,

\begin{equation}
	\mathbb{E}[\text{RIF}(Y; v)|X] = X \gamma + \epsilon,
\end{equation}

where the parameters $\gamma$ can be estimated by OLS.

In the case of quantiles, the influence function IF($Y, Q_\tau$) is given by $(\tau - \mathbb{I}\{Y \leq Q_\tau\}) / f_Y(Q_\tau)$, where $\mathbb{I}\{\cdot\}$ is an indicator function, $f_Y(\cdot)$ is the density of the marginal distribution of $Y$, and $Q_\tau$ is the population $\tau$-quantile of the unconditional distribution of $Y$. As a result, RIF($Y; Q_\tau$) is equal to $Q_\tau + \text{IF}(Y, Q_\tau)$, and can be rewritten as

\begin{equation}
	\text{RIF}(Y; Q_\tau) = Q_\tau + \frac{\tau - \mathbb{I}\{Y \leq Q_\tau\}}{f_Y(Q_\tau)} = c_{1,\tau} \cdot \mathbb{I}\{Y > Q_\tau\} + c_{2,\tau},
\end{equation}

where $c_{1,\tau} = \frac{1}{f_Y(Q_\tau)}$ and $c_{2,\tau} = Q_\tau - c_{1,\tau} \cdot (1 - \tau)$. Except for the constants $c_{1,\tau}$ and $c_{2,\tau}$, the RIF for a quantile is simply an indicator variable $\mathbb{I}\{Y \leq Q_\tau\}$ for whether the outcome variable is smaller or equal to the quantile $Q_\tau$. Using the terminology introduced above, running a linear regression of $\mathbb{I}\{Y \leq Q_\tau\}$ on $X$ is a distributional regression estimated at $y = Q_\tau$, using the influence function of the linear probability model.

There is, thus, a close connection between RIF regressions and the distributional regressional methods of 

Consider Firpo, Fortin, and Lemieux (2009) introduced the use of a specific type of IF, known as recentered influence functions (RIFs), useful for analyzing how changes in the distribution of explanatory variables (X) impact the unconditional distribution of the outcome variable (Y). {Rios-Avila2020stata}
\[
RIF(y; \nu, F_Y) = \nu(F_Y) + IF(y; \nu, F_Y)
\]
and demonstrated how the RIF-OLS could be used to estimate \( \nu(Y) \) as a series of regressions, applying to the distributional characteristic of interest, such as the quantiles. 

\section{RIF Regressions at Quantiles}

The RIF-OLS regression model allows us to estimate the effect of explanatory variables, $X$, on the unconditional quantile, $Q_{\tau}$, of an outcome variable, $Y$. The RIF is estimated in quantile regressions by first calculating the sample quantile $Q_{\tau}$ and computing the density at $Q_{\tau}$, that is $f(Q_{\tau})$ using kernel methods  {Firpo et al., 2009b}. Moreover, this approach relies on the indicator function $1\{Y \le Q_{\tau}\}$ taking value one if the condition in $\{\cdot\}$ is true, zero otherwise. Estimates for each observation $i$ of the RIF, $RIF(Y_{i}; Q_{\tau})$, are then obtained by inserting $Q_{\tau}$ and $f(Q_{\tau})$ in the aggregate RIF-function, defined as:

\begin{linenomath*}
	\begin{equation}
		\begin{split}
			RIF(Y; Q_{\tau}) &= Q_{\tau} + IF(Y; Q_{\tau}) \\
			&= Q_{\tau} + \frac{1 - 1\{Y \le Q_{\tau}\}}{f_{Y}(Q_{\tau})} \\
			&= \frac{1}{f_{Y}(Q_{\tau})}1\{Y > Q_{\tau}\} + Q_{\tau} - \frac{1}{f_{Y}(Q_{\tau})}(1 - \tau)
		\end{split}
	\end{equation}
\end{linenomath*}

where the RIF is the first order approximation of the quantile $Q_{\tau}$. $IF(Y; Q_{\tau})$ represents the influence function for the $\tau$th quantile. It measures the influence of an individual observation on the $\tau$th quantile. Adding the quantile $Q_{\tau}$ to the influence function yields the RIF. The probability density of $Y$ evaluated at $Q_{\tau}$ is $f_{Y}(Q_{\tau})$.

Firpo et al. (2009b) model the conditional expectation of the RIF-regression function, $\mathbb{E}[RIF(Y; Q_{\tau})|X]$, as a function of explanatory variables, $X$, in the UQR:

\begin{linenomath*}
	\begin{equation}
		\mathbb{E}[RIF(Y; Q_{\tau})|X] = g_{\tau}(X)
	\end{equation}
\end{linenomath*}

where a linear function $X \beta_{\tau}$ is specified for $g_{\tau}(X)$, as for example in Borah and Basu (2013). The average derivative of the UQR, $\mathbb{E}_{X}\left[ \frac{\partial g_{\tau}(X)}{\partial X} \right]$, captures the marginal effect of a small location shift in the distribution of covariates on the $\tau$th unconditional quantile of $Y$ keeping everything else constant. Therefore, the coefficients, $\beta_{\tau}$, can be unconditionally interpreted, as $\mathbb{E}[RIF(Y; Q_{\tau})] = \mathbb{E}_{X}[ \mathbb{E}[RIF(Y; Q_{\tau})|X]] = \mathbb{E}(X) \beta_{\tau}$. That is the unconditional expectations of $X$ in the unconditional quantile $Q_{\tau}$. This is one pitfall of CQRs in decomposition methods. The unconditional mean interpretation is important for decompositions in the sense of Oaxaca (1973) and Blinder (1973). Indeed, Oaxaca-Blinder decompositions use the unconditional mean interpretation of $\beta_{\tau}$ i.e. the interpretation of $\beta_{\tau}$ as the effect of increasing the mean value of $X$ on the mean value of $Y$. In UQR, the coefficients $\beta_{\tau}$ can thus be estimated by OLS in the following way:

\begin{linenomath*}
	\begin{equation}
		Q_{\tau} = \mathbb{E}[RIF(Y; Q_{\tau})|X] = \mathbb{E}(X) \beta_{\tau}
	\end{equation}
\end{linenomath*}

The basic wage equation of the RIF-OLS model at quantile $\tau$, with $\tau \in (0,1)$, is then:

\begin{linenomath*}
	\begin{equation}
		RIF(Y; Q_{\tau}) = X \beta_{\tau} + u_{\tau}
	\end{equation}
\end{linenomath*}

where $Y$ is the natural logarithm of hourly earnings and $X$ is a vector of $K$ explanatory variables (including the constant), $\beta_{\tau}$ is the corresponding coefficient vector and $u_{\tau}$ is the corresponding error term. The coefficient vector of the unconditional quantile at each observation $i$ is defined as:

\begin{linenomath*}
	\begin{equation}
		\beta_{\tau} = \left( \sum_{i=1}^{N} X_{i} X_{i}^{\prime} \right)^{-1} \sum_{i=1}^{N} X_{i} RIF(Y_{i}; Q_{\tau})
	\end{equation}
\end{linenomath*}

UQRs estimate the effect of covariates on all parts of the earnings distribution and are thus particularly interesting for policy implications or evaluation. CQRs do not allow to draw conclusions about the impact of a variable on the overall earnings distribution but rather provide insights about the dispersion of earnings within different subgroups of the target population {Borah and Basu, 2013}.

\subsection{Decomposition}

Given the assumptions that the mean of the RIF-function is equal to the actual quantile as well as to the mean of the conditional expectation given $X$ shown above, we have:

\begin{linenomath*}
	\begin{equation}
		\begin{split}
			\mathbb{E}[RIF(Y; Q_{\tau})|X]^{M} - \mathbb{E}[RIF(Y; Q_{\tau})|X]^{F} &= X_{m} \beta_{M} - X_{f} \beta_{F} \\
			&= \Delta_{\tau}
		\end{split}
	\end{equation}
\end{linenomath*}

where $\Delta_{\tau}$ is the GPG at the $\tau$th quantile and $M = Male$ and $F = Female$. The GPG is, as in the standard two-fold Oaxaca-Blinder decomposition, decomposed in an endowments (explained) and a coefficients (unexplained) component. The decomposition has
